% Deutsche Zusammenfassung; inhaltlich synchron mit preface/abstract.tex.

Code-Reviews sind auf die verfügbaren Informationen beschränkt. Verursacht
eine Änderung einen Fehler in einer anderen Datei, lässt sich diese Folge weder
aus dem Diff noch zuverlässig aus dem Gedächtnis ableiten. Diese Arbeit
untersucht, wie Repository-Kontext die inhaltliche Abdeckung LLM-generierter
Reviews verändert, ob struktureller Kontext dateiübergreifendes Schlussfolgern
ermöglicht und ob menschliche Bewertende den daraus resultierenden
Unterschied wahrnehmen.

Eine Pipeline erzeugt Reviews in vier Modi: eine reine Diff-Basislinie, ein
struktureller Arm, der Kanten aus einem mit Joern erstellten Code Property Graph
liefert, ein Retrieval-Arm, der embedding-ähnlichen Code liefert, und eine
Kombination beider Kontexte. Die Modi werden auf 35~Pull Requests aus fünf
Repositories, auf 40~kontrollierte Defektinjektionen und in einer
Nutzerstudie evaluiert.

Auf den Pull Requests bewertete ein Panel aus fünf Sprachmodellen jedes Review
anhand von 25~Kriterien. Der Graphmodus erzielt einen Zugewinn von
$+0{,}69$~Punkten bei den neun im Voraus als strukturell relevant festgelegten
Kriterien, und 96\,\% seines gesamten Zugewinns entfallen auf diese
Kriterien (jeweils $p = 0{,}003$). Retrieval erzielt den größten Zugewinn auf der
vollständigen 25-Kriterien-Skala ($+1{,}03$, $p = 0{,}002$). Nach der Korrektur
für multiple Vergleiche bleibt beim Graphmodus der gezielte Effekt und bei
Retrieval der Effekt auf der Gesamtskala signifikant; beim Hybridmodus bleibt
keiner der beiden Effekte signifikant.

Bei 28~injizierten Defekten, deren Folgen außerhalb der editierten Datei
liegen, erkennt die Basislinie keinen Defekt, Retrieval erkennt 5 und der
Graph~18. Das Hinzufügen von Vererbungskanten erhöht die Erkennung auf 26 von
28. Bei
12~Kontrolldefekten, die in der editierten Datei sichtbar sind, erreichen alle
Modi eine vollständige oder nahezu vollständige Erkennung. Die acht zusätzlich
durch Vererbungskanten erkannten Fälle zeigen erhebliches Verbesserungspotenzial
in der Graphkonstruktion.

In der Nutzerstudie verglichen 27~Teilnehmende Reviews mit Graph-Kontext
und Reviews der Basislinie. Beim Benennen betroffener Codestellen
bevorzugten sie das Review mit Graph-Kontext (mittlerer Präferenzwert
$0{,}710$; $0{,}5$ bezeichnet keine Präferenz; $p = 0{,}0005$). Bei der
wahrgenommenen Gesamtnützlichkeit zeigte sich keine Präferenz; bei den
Anmerkungen zur Codeklarheit bevorzugten sie das Review der Basislinie.

Der Beitrag besteht in einem klar abgegrenzten Fähigkeitsgewinn und nicht in
einer allgemeinen Verbesserung der Reviewqualität. Struktureller Kontext macht
dateiübergreifende Folgen sichtbar, die im Diff nicht enthalten sind.
Ähnlichkeitsbasiertes Retrieval erhöht dagegen die Abdeckung des gesamten
Kriterienrasters, ersetzt strukturelle Informationen jedoch nicht. Der Nutzen
ist groß, wenn dateiübergreifendes Schlussfolgern erforderlich ist. Im
Pull-Request-Benchmark fällt er kleiner, aber gezielt aus.
